\section{Matsubara \texorpdfstring{\acrlongpl{gf}}{GF}}\label{chap:gf}
In this section, we briefly motivate the idea behind \acrlongpl{gf} in the context of quantum many-body systems and introduce the machinery required for computing them at finite temperature via perturbation theory.
%\subsection{TODO Response function and self-energy}
%As alluded to in the introduction of this thesis, in condensed matter physics one already has a rather complete microscopic description in terms of the full Schrödinger equation for the many-body wavefunction \cite[p.5]{QMB_Coleman_2015}.
%One way to obtain information about a system is to measure how a perturbation evolves over time  .
\subsection{\texorpdfstring{\acrlongpl{gf}}{GF} in many-body physics}\label{sec:gf:motivation}
There are various extensive introductions to and discussions of the concept of \glspl{gf} for many-body physics in the literature, e.g. \cite{QFT_Abrikosov_1975,QMB_Mahan_2000,QMB_Bruus_2004,QMB_Fabrizio_2022}. %\cites[\schap 7]{QFT_Abrikosov_1975}[\schap 2 and 3]{QMB_Mahan_2000}[\schap 8]{QMB_Bruus_2004}[\schap 4]{QMB_Fabrizio_2022}
As an illustrative example, consider the function $\sgf_{a,b}(t)=-\i \bra{\psi_a(t) \psi_b^\dagger(0)}$ where $t>0$ and $a,b$ characterize the state created by the field operator $\psi^\dagger$, e.g. band indices, momenta or spin.
To be more explicit, at zero temperature the expectation value should be taken with respect to the ground state.
If $\psi_b^\dagger \Ket{0}$ is an eigenstate of the system, the time evolution merely contributes a phase factor and we only get a finite result for $a=b$.
Generally, however, this will not be an eigenstate and then the case $a=b$ contains information about how much of an excitation $b$ survives after evolving for time $t$.
The full object then also tells us how much of $b$ is scattered into different states \cite[\ssec 2.3]{QMB_Mahan_2000}.\\
\glspl{gf}, or correlation functions, enable us to calculate microscopic properties such as particle or current densities, and they contain information about a system's excitation spectrum \cite[\ssec 7.1]{QFT_Abrikosov_1975}.
From a pragmatic standpoint, being able to calculate \glspl{gf} can be viewed as a necessary step in evaluating expressions such as the Kubo formula, which measures the linear response of a system to an external perturbation \cite{QMB_Mahan_2000,QMB_Bruus_2004}.
Note that beyond the two-point functions involving only two operators, there exists an entire hierarchy of $N$-point functions \cite{AltlandSimons} and the concept of correlation functions also extends well beyond the notion of particle propagation.
For example, applying Kubo's formula to the evaluation of conductivity requires a correlation function involving current-density operators \cite{QMB_Coleman_2015}.\\
Suffice it to say that \glspl{gf} are ubiquitous and we need to be able to calculate them.
The situation is slightly more complicated than suggested by the example above once we include temperature, as is discussed in the upcoming section.
Since we always work in equilibrium here, we choose the Matsubara formalism to compute \glspl{gf}.
Bowen claims that it is the ``simplest finite temperature formalism'' \cite[\spage 5]{Bosonization_FiniteT_Bowen_2001}, and \cite{Giamarchi_2004,Gogolin_1999,Fradkin_2013,QMB_Fabrizio_2022,Bosonization_TLL_Schonhammer_1997,Bosonization_GF_LL_Meden_1999,Bosonization_Rao_2001,Bosonization_LowE_1d_fluids_Haldane_1981} all making use of imaginary time suggests that the corresponding authors came to a similar conclusion.
It should at least be noted that, in principle, one could instead also use the Keldysh approach \cite[\schap 11]{AltlandSimons}.
The latter is designed to also work outside equilibrium, so unsurprisingly, it is quite a bit more involved \cite{GF_intro_InOutEquil_Ge_2023}.
However, one advantage over the Matsubara method is that Keldysh works directly in real time, and thus avoids an analytic continuation from imaginary to real frequencies further down the line.
Technicalities aside, this continuation can in the end boil down to little more than a simple replacement of symbols, but of course this requires obtaining an intermediate result that is analytic in the relevant quantities.
Numerically the continuation appears to be rather difficult \cite{NumericAnalyticCont_QMC_MaxEntropy_Gubernatis_1991,NumericAnalyticCont_ImagTimeData_Burnier_2011}, so within the confines of this thesis we try to avoid this by working out analytic expressions for our \glspl{gf}.
\subsection{Interaction picture and imaginary time}\label{sec:gf:interaction_imaginary_time}
Here we mostly follow \cite{QMB_Bruus_2004}, although the presentation in \cites[\schap 4]{QMB_Fabrizio_2022}[\schap 5]{QMB_Nayak_2004} is very similar.
We always work in the interaction picture defined by the equations \cite[\ssec 5.3]{QMB_Bruus_2004}
\begin{align}
\Ket{\psi_\tint(t)} &= U_0(t)^\dagger \Ket{\psi(t)} \quad \text{and} \quad A_\tint(t) = U_0(t)^\dagger A  U_0(t)
\end{align}
where the full Hamiltonian $H=H_0+V$ is split into a free part $H_0$ and a perturbation $V$.
The operator $U_0(t) = \powe{-\i H_0 t}$ corresponds to the time evolution with respect to the free Hamiltonian and is assumed to be sufficiently simple to allow the computation of eigenvalues and eigenstates.
Generally speaking, this implies that $H_0$ is quadratic in the relevant degrees of freedom, e.g. a kinetic energy term.
Note that $H$ must be time-independent since we assume to be in equilibrium, and thus the full time evolution operator takes the simple form $U(t,t') = \powe{-\i H(t-t')}$.
Also, since we work in the grand canonical ensemble, the Hamiltonian here is really $H - \mu N$ \cite{QFT_FiniteT_Kapusta_2023,QMB_Coleman_2015,GF_SelfEnergyFunctionals_Eder_2021}, but as seems to be customary, we still only write $H$ for convenience \cite[\spage 27]{QMB_Bruus_2004}. % QFT FiniteT Kapusta p.2, Eder GF p.3, Coleman [search "grand"], Bruus\\
Inserting the definition of the state in the interaction picture into the Schrödinger equation yields
\begin{align*}
\i\partial_t \Ket{\psi(t)} &= \i \partial_t U_0(t) \Ket{\psi_\tint(t)} = H_0 U_0(t) \Ket{\psi_\tint(t)} + U_0(t) \i\partial_t \Ket{\psi_\tint(t)} \overset{!}{=} H \Ket{\psi(t)},
\end{align*}
and cancelling the free part on both sides of the equation then leads to
\begin{align}
\i\partial_t \Ket{\psi_\tint(t)} &= U_0(t)^\dagger V(t) \Ket{\psi(t)} = V_\tint(t) \Ket{\psi_\tint(t)}. \label{eqn:gf:interaction_picture_schrodinger_vint}
\end{align}
Writing $\Ket{\psi_\tint(t)} = U_\tint(t,t_0) \Ket{\psi_\tint(t_0)}$ where $\Ket{\psi_\tint(t_0)}$ is some arbitrary initial state in the interaction picture, we find that the same equation must also be satisfied by $U_\tint(t,t_0)$ itself.
Thus the full time evolution, given by $U(t,t') = U_0(t) U_\tint(t,t') U_0(t')^\dagger$, only requires us to solve \autoref{eqn:gf:interaction_picture_schrodinger_vint}.
At least formally, this is possible by iteratively integrating the equation, ultimately leading to \cite[\spage 90]{QMB_Bruus_2004}
\begin{align*}
U_\tint(t,t') &= \mathcal{T}_t \exp \kla*{-\i \intx{t'}{t}{V_\tint(t_1)}{t_1}} \equiv \sum_{n\ge 0} \frac{(-\i)^n}{n!} U_\tint^{(n)}(t,t')
\end{align*}
where the time-ordering operator $\mathcal{T}_t$ is defined by
\begin{align}
\mathcal{T}_t A(t_1) B(t_2) &= \begin{cases} A(t_1) B(t_2) & \text{for } t_1\ge t_2,\\ sB(t_2) A(t_1) & \text{for } t_1<t_2. \end{cases} \label{eqn:gf:time_order_def}
\end{align}
Here, $s=-1$ if both $A$ and $B$ are fermionic operators and $s=+1$ otherwise \cites[\spage 122]{QMB_Fabrizio_2022}. %[\spage 98]{QMB_Coleman_2015} \cites[\spage 160]{QMB_Bruus_2004}. %https://physics.stackexchange.com/questions/194073/minus-sign-in-the-time-ordering-operator
It is rather straight-forward to check that this is indeed a solution of \autoref{eqn:gf:interaction_picture_schrodinger_vint} by explicitly differentiating
\begin{align*}
U_\tint^{(n)}(t,t') &= \intx{t'}{t}{\dots \intx{t'}{t}{ \mathcal{T}_t V_\tint(t_1) \dots V_\tint(t_n)}{t_n} \dots }{t_1}
\end{align*}
with respect to $t$.
The product rule generates $n$ terms, and in each of them the Leibniz rule reduces one of the integrals to its integrand evaluated at $t$, i.e. $V_\tint(t)$.
Since all of the integration variables are smaller than $t$, the time-ordering then allows us to pull out this term to the left.
Up to relabelling the terms are all identical, and thus $\partial_t U_\tint^{(n)}(t,t') = n V_\tint(t) U_\tint^{(n-1)}(t,t')$, which directly leads to $\partial_{t} U_\tint(t,t') = -\i V_\tint(t) U_\tint(t,t')$.\\
The key difference between zero and finite temperature with respect to the computation of \glspl{gf} lies in the way we evaluate expectation values \cite{QMB_Mahan_2000}.
For some operators $A$ and $B$, we might consider the time-ordered or causal correlation function $C^\c_{AB} = \bra*{\mathcal{T}_t A(t) B(t')}$, where the expectation value at finite temperate is given by an ensemble average over all possible states, i.e.
$
\bra*{A} = \frac{1}{Z} \tr \klc*{ \powe{-\beta H} A}
$
with the partition function $Z=\tr \powe{-\beta H}$.
Using the full time evolution operator $U$, we can rewrite this correlation function as \cite[\schap 10]{QMB_Bruus_2004}
\begin{align*}
C^\c_{AB}(t,t') &= \frac{1}{Z} \tr \klc*{ \powe{-\beta H} \mathcal{T}_t U(t,0)^\dagger A U(t,0) U(t',0)^\dagger B U(t',0)} \\
&= \frac{1}{Z} \tr \klc*{ \powe{-\beta H} \mathcal{T}_t U_\tint(t,0)^\dagger A_\tint(t) U_\tint(t,0) U_\tint(t',0)^\dagger B_\tint(t') U_\tint(t',0)}.
\end{align*}
At this point, the commonly cited reason for the need of a new formalism is that we would now have to expand both the density matrix and the time evolution operator in powers of $V$, leading to a significantly more complicated perturbation series \cites[\schap 3]{QMB_Mahan_2000}. % [Mahan ch.3, BF ch.10]
By going to imaginary time we can circumvent this problem, because both of these terms may then be treated in one go.\\
The Matsubara correlation function is defined as
\begin{align}
C^\M_{AB}(\tau,\tau') &= -\bra*{\mathcal{T}_\tau A(\tau) B(\tau')}
\end{align}
where $\tau=\i t$ is referred to as the \emph{imaginary time} and $\mathcal{T}_\tau$ is the corresponding time-ordering operator.
This should merely be seen as a mathematical trick that allows us to write down a perturbation series analogous to the zero temperature case; from a physical point of view, both time and frequency must be real-valued quantities \cite[\spage 158]{QMB_Bruus_2004}.
In \autoref{sec:gf:Kallen_Lehmann}, we discuss how the -- physically relevant -- retarded correlation function can be obtained from the Matsubara one, but first let us summarize the formalism.\\
Time evolution is no longer unitary, thus $U_0(\tau)^{-1} = \powe{+\tau H_0}$ is the inverse to $U_0(\tau)=\powe{-\tau H_0}$.
The rules for transforming to the interaction picture remain identical, e.g. $A_\tint(\tau) = U_0(\tau)^{-1} A U_0(\tau)$.
The full time evolution operator becomes $U(\tau,\tau') = U_0(\tau) U_\tint(\tau,\tau') U_0(\tau')^{-1}$, where the interacting part is given by
\begin{align*}
U_\tint(\tau,\tau') &= \mathcal{T}_\tau \exp \kla*{- \intx{\tau'}{\tau}{V_\tint(\tau_1)}{\tau_1}}.
\end{align*}
Note that the density operator can be represented using the time evolution operator via
\begin{align*}
\powe{-\beta H} &= U(\beta,0) = \powe{-\beta H_0} U_\tint(\beta,0)
\end{align*}
as desired.
The Matsubara correlation function can then be rewritten as
\begin{align*}
C^\M_{AB}(\tau,\tau') &= -\frac{1}{Z} \tr \klc*{ \powe{-\beta H} \mathcal{T}_\tau U(0,\tau) A U(\tau,0) U(0,\tau') B U(\tau',0) }\\
%&= -\frac{1}{Z} \tr \klc*{ \powe{-\beta H_0} U_\tint(\beta,0) \mathcal{T}_\tau U^\dagger(\tau,0) A U(\tau,0) U^\dagger(\tau',0) B U(\tau',0) }\\
&= -\frac{1}{Z} \tr \klc*{ \powe{-\beta H_0} U_\tint(\beta,0) \mathcal{T}_\tau U_\tint(0,\tau) A_\tint(\tau) U_\tint(\tau,\tau') B_\tint(\tau') U_\tint(\tau',0)} \\
&= -\frac{1}{Z} \tr \klc*{\powe{-\beta H_0} \mathcal{T}_\tau U_\tint(\beta,0) A_\tint(\tau) B_\tint(\tau')}.
\end{align*}
The last equality results from exploiting the time-ordering to rearrange terms such that $U_\tint(\beta,\tau) U_\tint(\tau,\tau') U_\tint(\tau',0) = U_\tint(\beta,0)$ \cite[\spage 135]{QMB_Fabrizio_2022}.
Since the partition function can also be expressed as an expectation value with respect to the free Hamiltonian via
\begin{align*}
Z &= \tr \powe{-\beta H} = \tr \klc*{\powe{-\beta H_0} U_\tint(\beta,0)} = Z_0 \bra*{U_\tint(\beta,0)}_0,
\end{align*}
we can finally rewrite the full correlation function as
\begin{align}
C^\M_{AB}(\tau,\tau') &= -\frac{\bra*{\mathcal{T}_\tau U_\tint(\beta,0) A_\tint(\tau) B_\tint(\tau')}_0}{\bra*{U_\tint(\beta,0)}_0}. \label{eqn:gf:correlator_Matsubara_AB_bra0_int}
\end{align}
The structure here is remarkably similar to the one for zero temperature, and this result readily generalizes to correlators with an arbitrary number of operators \cite[\seqn 4.40]{QMB_Fabrizio_2022}.\\
Using the cyclic property of the trace, one can show that the correlation function only depends on the difference between the time arguments, and that it is (anti-) periodic for bosons (fermions) with frequency $\beta$ \cite[\ssec 10.1]{QMB_Bruus_2004}, i.e.
\begin{align*}
C^\M_{AB}(\tau,\tau') &= C^\M_{AB}(\tau-\tau') \quad \text{and} \quad C^\M_{AB}(\tau+\beta) = \pm C^\M_{AB}(\tau).
\end{align*}
For any periodic function satisfying $f(\tau+\beta) = \pm f(\tau)$ we can write down a \gls{ft} with discrete frequencies via \cite[\ssec 5.1]{QMB_Nayak_2004}
\begin{align*}
f(\tau) &= \frac{1}{\beta} \sum_{n\in\mathbb{Z}} f_n \powe{-\i \omega_n \tau}.
\end{align*}
Assuming $\powe{\i (\omega_n - \omega_m)\beta} = 1$, we can compute the coefficients by inserting a resolution of unity
\begin{align*}
f_n &= \sum_{m\in\mathbb{Z}} f_m \delta_{n,m} = \sum_{m\in\mathbb{Z}} f_m \frac{1}{\beta} \intx{0}{\beta}{\powe{\i (\omega_n - \omega_m)\tau}}{\tau} = \intx{0}{\beta}{f(\tau) \powe{\i\omega_n \tau}}{\tau}
\end{align*}
which also defines the inverse of the transform.
The (anti-)periodicity then requires $\powe{\i \omega_n \beta} = \pm 1$ and leads to the definition of the bosonic and fermionic Matsubara frequencies,
\begin{align}
\omega_n^\B &= \frac{2\pi n}{\beta} \quad\text{and}\quad \omega_n^\F = \frac{(2n+1) \pi}{\beta}, \label{eqn:gf:Matsubara_freq_def}
\end{align}
that both satisfy the assumption $\powe{\i (\omega_n - \omega_m)\beta} = 1$.
Note that by convention one usually writes $f_n \equiv f(\i\omega_n)$ for the Fourier coefficients in order to simplify the notation for the generalization, or analytic continuation, of such functions to arbitrary complex frequencies \cite{QMB_Bruus_2004}.
\subsection{Perturbation series}\label{sec:gf:perturbation_series}
To actually evaluate \autoref{eqn:gf:correlator_Matsubara_AB_bra0_int}, one can write down a perturbation series in the interaction.
The expansion of the time-evolution operator is
\begin{align*}
U_\tint(\beta, 0) &= 1 - g\intx{0}{\beta}{V_\tint(\tau)}{\tau} + \frac{g^2}{2} \intx{0}{\beta}{\intx{0}{\beta}{\mathcal{T}_\tau V_\tint(\tau) V_\tint(\tau')}{\tau'} }{\tau} + \dots \numberthis\label{eqn:gf:U_int_expansion}
\end{align*}
and, in the following, for convenience of notation the time-ordering is implicit.
With the action $\sact_\tint = \intx{0}{\beta}{V_\tint}{\tau}$ and $\mathcal{O} = A_\tint(\tau) B_\tint(\tau')$ we then have
\begin{align*}
C^\M_{AB}(\tau,\tau') &= - \frac{\bra*{\powe{-\sact_\tint} \mathcal{O} }}{\bra*{\powe{-\sact_\tint}}} = - \frac{\bra*{\mathcal{O} \kla*{1 - \sact_\tint + \frac{1}{2} \sact_\tint^2 + \dots}}}{1 - \bra*{\sact_\tint} + \frac{1}{2} \bra*{\sact_\tint^2} + \dots}\\
&= -\bra[\Big]{\mathcal{O} \kla[\Big]{1 - \sact_\tint + \frac{1}{2} \sact_\tint^2 + \dots}} \klb[\Big]{1 + \bra*{\sact_\tint} - \frac{1}{2} \bra*{\sact_\tint^2} + \bra*{\sact_\tint}^2 + \dots}\\
&= -\bra*{\mathcal{O}} + \bra*{\mathcal{O} \sact_\tint}^\c - \frac{1}{2} \bra*{\mathcal{O};\sact_\tint^2}^\c - \bra*{\mathcal{O}; \sact_\tint}^\c \bra*{\sact_\tint} + \dots \numberthis\label{eqn:gf:perturbation_series}
\end{align*}
where $\bra*{A;B}^\c = \bra*{AB} - \bra*{A}\bra*{B}$ denotes the connected average\footnote{We include the semicolon when there are more than two terms to avoid ambiguity.} \cites[\ssec 11.3]{QFT_RG_Fradkin_2021}[\ssec 3.7.2]{QMB_Refresher_Tremblay_2008}.
Usually one would now think about how the combinatoric explosion of this expansion can be handled via diagrammatic techniques and Feynman rules \cite{QMB_Bruus_2004,QMB_Coleman_2015}.  %maybe Coleman or Bruus
In our case the calculation becomes practically infeasible well before combinatorics hinder progress, so we have no need for such advanced representations.
Therefore, the first order of \autoref{eqn:gf:perturbation_series} is already sufficient for our calculation in \autoref{sec:ell:gf}.
\subsection{\nameKallen--Lehmann representation}\label{sec:gf:Kallen_Lehmann}
Let us start by assuming to have a complete set of eigenstates and eigenvalues for the full Hamiltonian, i.e. $H\Ket{n} = E_n \Ket{n}$.
Of course if this was possible, we would not have any need for approximation schemes, and, in all likelihood, be dealing with a rather simple system.
Formally, however, we may then write down a rather simple representation of the correlation function by explicitly evaluating the trace.
This is possible by inserting two resolutions of unity to get \cite[\ssec 10.2]{QMB_Bruus_2004}
\begin{align*}
C^\M_{AB}(\tau) &= -\bra*{\mathcal{T}_\tau A(\tau) B(0)} = -\frac{1}{Z} \tr \klc*{\powe{-\beta H} \powe{\tau H} A \powe{-\tau H} B}\\
&= -\frac{1}{Z} \sum_{n,n'} \brc[\big]{n}{ \powe{-\beta H} \powe{\tau H} A }{n'} \brc[\big]{n'}{ \powe{-\tau H} B}{n}\\
&= -\frac{1}{Z} \sum_{n,n'} \powe{-\beta E_n} \powe{\tau (E_n - E_{n'})} \brc{n}{A}{n'} \brc{n'}{B}{n}.
\end{align*}
In frequency space this gives rise to the \emph{\nameKallen--Lehmann}\footnote{Although except for \cite{QMB_Coleman_2015} the other references given in this section only seem to credit Lehmann.} representation \cite[\spage 162]{QMB_Bruus_2004}
\begin{align*}
C^\M_{AB}(\i\omega_n) &= \intx{0}{\beta}{\powe{\i\omega_n\tau} C^\M_{AB}(\tau) }{\tau} = -\frac{1}{Z} \sum_{n,n'} \powe{-\beta E_n} \brc{n}{A}{n'} \brc{n'}{B}{n} \intx{0}{\beta}{\powe{\i\omega_n\tau} \powe{\tau (E_n - E_{n'})} }{\tau}\\
%&= -\frac{1}{Z} \sum_{n,n'} \powe{-\beta E_n} \frac{\brc{n}{A}{n'} \brc{n'}{B}{n}}{\i\omega_n + E_n - E_{n'}} \kla*{\powe{(\i\omega_n + E_n-E_{n'}) \beta} - 1 }.\\
&= \frac{1}{Z} \sum_{n,n'} \frac{\brc{n}{A}{n'} \brc{n'}{B}{n}}{\i\omega_n + E_n - E_{n'}} \kla*{\powe{-\beta E_n} \mp  \powe{-\beta E_{n'}}},
\end{align*}
where as usual the upper sign is for bosons and the lower sign for fermions.
The crucial observation is that now, having thrown away any concerns about how to obtain these eigenstates in practise, we can also directly compute the correlation function in real time.
Similar to the above one finds  \cite[\spage 134]{QMB_Bruus_2004}
\begin{align*}
C^\ret_{AB}(t) &= \i\Theta(t) \bra*{\klb*{A(t), B(0)}_\mp} = \i\Theta(t) \frac{1}{Z} \sum_{n} \powe{-\beta E_n} \brc[\big]{n}{\klb*{\powe{\i Ht} A \powe{-\i Ht}, B}_\mp }{n}\\
&= \i\Theta(t) \frac{1}{Z} \sum_{n,n'} \powe{-\beta E_n} \kla*{\brc{n}{\powe{\i Ht} A \powe{-\i Ht} }{n'} \brc{n'}{B}{n} \mp \brc{n}{B}{n'}\brc{n'}{\powe{\i Ht} A \powe{-\i Ht}}{n}}\\
&= \i\Theta(t) \frac{1}{Z} \sum_{n,n'} \brc{n}{A}{n'} \brc{n'}{B}{n} \powe{\i (E_n - E_{n'}) t} \kla*{\powe{-\beta E_n} \mp \powe{-\beta E_{n'}}},
\end{align*}
where the summation indices are swapped in the second line.
In order to make the \gls{ft} converge we have to include an infinitesimal positive imaginary part in the frequency.
Then, using $\intx{0}{\infty}{\i \powe{\i (E + \i\eta) t} }{t} = \frac{1}{E + \i \eta}$ gives
\begin{align*}
C^\ret_{AB}(\omega + \i \eta) &= \intr{\powe{\i(\omega + \i \eta) t} C^\ret_{AB}(t)}{t} = \frac{1}{Z} \sum_{n,n'} \frac{\brc{n}{A}{n'} \brc{n'}{B}{n} }{\omega + \i\eta + E_n - E_{n'}} \kla*{\powe{-\beta E_n} \mp \powe{-\beta E_{n'}}}.
\end{align*}
Note that this is structurally identical to the result found for the Matsubara correlator, which suggests that ``the'' correlation function can be defined as \cite[\seqn 10.29]{QMB_Bruus_2004}
\begin{align*}
C_{AB}(\omega) &= \frac{1}{Z} \sum_{n,n'} \frac{\brc{n}{A}{n'} \brc{n'}{B}{n} }{\omega + E_n - E_{n'}} \kla*{\powe{-\beta E_n} \mp \powe{-\beta E_{n'}}}
\end{align*}
for any $\omega\in\mathbb{C}$.
Thus we can obtain the retarded \gls{gf} by means of an analytic continuation $\i\omega_n \rightarrow \omega + \i\eta$ in the Matsubara correlator. %derivation Bruus ~eq.10.25
Similarly, the advanced \gls{gf} is found by $\i\omega_n \rightarrow \omega - \i\eta$.\\
Importantly, an integral can generally not be considered an analytic expression.
For example, substituting $\i\omega_n = \omega + \i\eta$ in the phase factor $\powe{\i\omega_n \tau}$ in the derivation of the \nameKallen--Lehmann representation would give a different -- and incorrect -- result \cite[\spage 163]{QMB_Bruus_2004}.
%This only seems to be a problem for the integration over $\tau$ though